%% drtran -- Box-Jenkins transfer function models.
%% Copyright (C) 1995-2026 A.B. Treadway, J.A. Mauricio & D.E. Guerrero.
%%
%% This document is free: you may redistribute it and/or modify it under the
%% terms of the GNU General Public License as published by the Free Software
%% Foundation; either version 2 of the License, or (at your option) any later
%% version.  Distributed WITHOUT ANY WARRANTY; see COPYING for details.
%%
%% Build:  pdflatex drtran-note

\documentclass[11pt,a4paper]{article}

\usepackage[T1]{fontenc}
\usepackage[utf8]{inputenc}
\usepackage{lmodern}
\usepackage{amsmath}
\usepackage{amssymb}
\usepackage{booktabs}
\usepackage{geometry}
\usepackage{verbatim}
\usepackage[hidelinks]{hyperref}

\geometry{margin=2.7cm}

\title{\bfseries Re-implementing Box--Jenkins transfer function models\\
       on an exact VARMA likelihood:\\[2pt]
       \large engineering notes, and one identification hazard}

\author{D. E. Guerrero\thanks{Universidad Complutense de Madrid,
        \texttt{davidesg@ucm.es}. This note documents the \texttt{drtran}
        program, which is free software under the GNU GPL. The exact
        likelihood engine is J. A. Mauricio's and is used unmodified; the
        methodology is A. B. Treadway's school. Errors are mine.}}

\date{July 2026}

\begin{document}
\maketitle

\begin{abstract}
\noindent
\textbf{Nothing in this note is new.} The models are Box and Jenkins's, the
identification procedure is theirs, the exact multivariate likelihood is
Mauricio's, and the methodology is Treadway's. What the note reports is what
happens when that classical apparatus is \emph{re-assembled} on top of an exact
vector ARMA likelihood, in one program, with everything estimated at once:
a handful of engineering decisions turn out to depart from the original
solutions in ways that are measurable, and one of them exposes an
identification hazard that is easy to walk into and easy to miss. We show that
the contemporaneous transfer coefficient $\omega_0$ and the innovation
covariance $\sigma_{12}$ are \emph{exactly} observationally equivalent when the
noise and the input share the same autoregressive operator, and are separated
otherwise only by a signal proportional to the difference between those
operators. In a real case (Spanish CPI and the price of crude oil) that
separating signal is $10\%$ of the main one, and the likelihood develops a flat
ridge along which the fit does not improve while the estimates run to a corner.
Every number below is reproducible from the program's test battery.
\end{abstract}

\section{What is not new}
\label{sec:notnew}

It is worth being blunt about this first, because the rest of the note is only
interesting once it is settled.

\begin{itemize}
\item \textbf{The models} are the transfer function models of
  Box \& Jenkins (1970), chapters 10--11 of the current
  edition~\cite{bjr1994}.

\item \textbf{The cast} used here---write $N_t = Y_t - \sum_j \nu_j(B) X_{j,t}$,
  stack $N$ and the inputs into a vector, and note that the resulting system is
  a VARMA with \emph{diagonal} $\Phi$, $\Theta$ and $Q$---is the standard
  multivariate representation of a transfer function model. It is in Tiao and
  Box~\cite{tiaobox1981}, in Hillmer and Tiao, and in Reinsel. It is also
  written out by hand, for a three-variable system, in
  Muñoz Polo (2001)~\cite{munoz2001}, \S2.4.

\item \textbf{The network} of \S\ref{sec:network}---several outputs, a series
  that is at once the output of one relation and the input of another---is a
  \emph{recursive} dynamic simultaneous system: a Wold causal chain. It is what
  Mauricio's \texttt{m6} labour-market models already were in 1996; they were
  simply hand-coded, one \texttt{shootx} at a time.

\item \textbf{The scale normalization} $Q_{11}=1$ of \S\ref{sec:qq} is textbook
  identification under a concentrated scale factor. Mauricio himself states that
  the decomposition $\Sigma = \sigma^2 Q$ is not unique~\cite{mauricio1995,
  mauricio2002}.

\item \textbf{The exact Hessian} of \S\ref{sec:hessian} is ordinary numerical
  practice~\cite{dennis1983}. The code for it was already present in the
  ancestor programs---\emph{commented out}.

\item \textbf{The identification hazard} of \S\ref{sec:identification} is, at
  bottom, the Cowles Commission problem: one cannot free a contemporaneous
  structural coefficient and the innovation covariance at the same time. That is
  eighty years old and is the identification problem of any structural VAR. What
  \S\ref{sec:identification} adds is only a quantitative refinement in the
  transfer-function case, and a diagnostic that follows from it.
\end{itemize}

\noindent
What \emph{is} contributed here is engineering, not science: an automatic
reduction from a univariate model specification to a VARMA structure; a
declarative way of stating parameter restrictions that were previously written
into code by hand; three concrete software-correctness findings about a family
of programs that has been in use for thirty years; and a test battery that pins
the result to the univariate program's output to six decimal places. This note
is therefore a \emph{software and reproducibility} note. It is offered in that
spirit.

\section{The lineage, and the one structural idea}

\subsection{Lineage}

\begin{center}
\begin{tabular}{ll}
\toprule
Box \& Jenkins (1970) & the models \\
\quad Jenkins's program & \\
\quad\quad \textsc{taste} (Treadway) & identification and estimation; backforecasting \\
\quad\quad\quad \texttt{fue} & univariate (ARIMA/GSM with intervention) \\
Mauricio (1995, 1997) & \emph{exact} likelihood of a VARMA --- \texttt{elf} \\
\quad \texttt{drvarma} & estimates arbitrary VARMA models \\
\quad\quad \texttt{drtran} & the bridge: \texttt{fue} $\rightarrow$ VARMA $\rightarrow$ \texttt{elf} \\
\bottomrule
\end{tabular}
\end{center}

Muñoz Polo's dissertation~\cite{munoz2001}, supervised by Treadway, is the
document that fixes the methodology of the school, and it already estimates with
Mauricio's engine (``the software implementing the exact ML algorithm programmed
and designed by Mauricio (1995, 1996, 1997)'', ch.~3). The two halves of the
lineage were therefore already joined in 2001---\emph{by hand}. What
\texttt{drtran} does is automate the joining.

\subsection{The cast}

A transfer function model
\begin{equation}
Y_t \;=\; \sum_{j=1}^{k} \frac{\omega_j(B)}{\delta_j(B)}\,B^{b_j}\,X_{j,t}
        \;+\; N_t
\label{eq:tf}
\end{equation}
is recast as a VARMA of $m = 1+k$ series in which
\begin{align}
w_{t,1} &= w^{Y}_t - \textstyle\sum_j \nu_j(B)\, w^{X_j}_t \qquad (\text{the noise } N),\\
w_{t,j+1} &= w^{X_j}_t,
\end{align}
with $\Phi$ and $\Theta$ \emph{diagonal} (each series carries its own ARMA) and
$Q$ diagonal. All of the coupling lives in the \emph{subtracted} transfer, so the
VARMA that the likelihood engine sees has no off-diagonal element anywhere.

Two consequences follow, and they are the whole design:

\begin{enumerate}
\item Mauricio's engine can be used \textbf{unmodified}. In this project any
  discrepancy with the univariate program was treated, without exception, as a
  bug in the bridge. It always was one.
\item The joint fit \textbf{reproduces} the univariate program run separately on
  each series, to six decimals, on six real models
  (deterministic seasonality; airline; generalized seasonality with
  fixed-frequency factors; a fixed AR coefficient). This is the entry gate of the
  test battery, and \S\ref{sec:separability} explains why it must hold.
\end{enumerate}

\section{Six engineering departures, measured}

\subsection{The input model: fixed, or estimated jointly}
\label{sec:separability}

The doctrine is explicit~\cite[\S2.6]{munoz2001}: \emph{``the univariate model of
the input remains unaltered from the beginning to the end of the process.''}
\texttt{drtran} does the opposite: it estimates everything---the output's ARMA,
the input's ARMA, the transfer, the deterministic terms and the means---in a
single maximization.

\textbf{Does it matter? No---and that is the result.} On Spanish CPI
$\leftarrow$ WTI crude:

\begin{center}
\begin{tabular}{lrr}
\toprule
 & input ARMA \emph{fixed} & input ARMA \emph{free} \\
\midrule
$\log L$   & $-718.183933$ & $-718.183933$ \\
$\omega_0$ & $0.016402$    & $0.016402$ \\
$\omega_1$ & $0.010792$    & $0.010792$ \\
\bottomrule
\end{tabular}
\end{center}

Identical. This is not a coincidence. With $Q$ \emph{diagonal} and no feedback,
the joint likelihood \emph{factorizes} as $L(N)\cdot L(X)$, and $L(X)$ does not
involve the transfer. Box and Jenkins's decision to leave the input model alone
is therefore \textbf{not an approximation: it is exactly optimal} in that case.
It is also, precisely, why the joint fit homologates with the univariate program.

The factorization \textbf{breaks the moment $\sigma_{12}$ is freed}---and that is
exactly the case the dissertation~\cite[\S2.4]{munoz2001} points to as the
starting point for specifying a bivariate relation. There, joint estimation stops
being a luxury.

\subsection{$\sigma^2 Q$: the normalization that was missing}
\label{sec:qq}

The engine assumes $a_t \sim N(0,\sigma^2 Q)$ and concentrates $\sigma^2$ out.
Mauricio states in both references that this decomposition is \emph{not unique},
and that this \emph{``raises no problem in the estimation of $E[a_t a_t']$, since
interest lies in the product $\sigma^2 Q$, but not in $\sigma^2$ nor in $Q$ by
itself''}~\cite[p.~474]{mauricio2002}.

\textbf{He is right, and we verified it.} The concentrated criterion
is~\cite[eq.~12]{mauricio2002}
\begin{equation}
F(\beta\mid w) \;=\; \bigl(\tilde{w}' R_W^{-1}\tilde{w}\bigr)\,
                     |R_W|^{1/(mn)},
\label{eq:crit}
\end{equation}
and under $Q\to cQ$ the matrix $R_W$ scales by $c$: the first factor goes as
$1/c$, the second as $c$. \emph{$F$ is exactly invariant.} Since $\sigma^2$ is
concentrated, $\hat\sigma^2(cQ)\cdot cQ = \hat\sigma^2(Q)\cdot Q$: the product
does not move. Estimating the same model with the scale of $Q$ free (as the
ancestor programs do) and with it normalized:

\begin{center}
\begin{tabular}{lrr}
\toprule
 & scale free & $Q_{11}=1$ \\
\midrule
$\log L$                         & $-767.424341$ & $-767.424341$ \\
$\hat\Sigma$                     & $0.062666 \,/\, 68.838114$ & $0.062666 \,/\, 68.838114$ \\
$\mathrm{SE}(\phi_1)$            & $0.062204$ & $0.062204$ \\
$\mathrm{SE}(\mu)$               & $0.028502$ & $0.028502$ \\
$\mathrm{SE}(\log \mathrm{var}_2/\mathrm{var}_1)$ & $\mathbf{0.193083}$ & $\mathbf{0.136399}$ \\
\bottomrule
\end{tabular}
\end{center}

The flat direction is \emph{exact}, and it is harmless for $\log L$, for
$\hat\Sigma$, and for the standard errors of $\phi,\theta,\mu,\omega,\delta$. The
damage is confined to \textbf{inference on $Q$ itself}: the scale parameter,
which is \emph{literally} unidentified, receives a respectable-looking
$t=-0.19$ with $p=0.85$; and the variance ratio, which \emph{is} identified,
comes out with a standard error inflated by $42\%$.

\texttt{drtran} normalizes $Q_{11}=1$ for one concrete reason: \textbf{deciding
whether a covariance is zero is a test}, and a test has to be a real one. Nothing
else needed it.

\subsection{Standard errors: exact Hessian, not the accumulated one}
\label{sec:hessian}

The ancestor programs compute standard errors from the \emph{BFGS-accumulated}
Hessian. In all of them the exact alternative (a finite-difference Hessian plus a
Cholesky factorization) is written and \emph{commented out}. It was a 1996 cost
decision that nobody revisited.

A BFGS Hessian is built from the steps the optimizer \emph{takes}; it therefore
never sees the directions in which the optimizer does not move, and produces
finite numbers where there should be none. Two measured consequences: the
univariate program's standard errors are \emph{unstable}---two runs of the same
model, with identical point estimates, report $\mathrm{SE}(\mu)=0.073304$ and
$0.028316$; and the standard errors reported for the nine elements of $Q$ in
\texttt{m6-1} lie along the flat direction of \S\ref{sec:qq}.

Using the exact Hessian, $\mathrm{SE}(\mu)=0.028502$, which is the exact GLS
value.

\subsection{``A single output'' versus the network}
\label{sec:network}

The title of \S2.5 of the dissertation is, literally, \emph{``Representation of
the Single-Output Transfer Model''}. That is the boundary of the classical
apparatus: one output, $k$ exogenous inputs.

The school's real systems are not that. In \texttt{m6} (the Spanish labour
market), \texttt{EU} is an \emph{output} of \texttt{EC} \emph{and} an \emph{input}
to \texttt{EI}:
\begin{center}
\texttt{EC} $\xrightarrow{\;b=2\;}$ \texttt{EU} $\xrightarrow{\;b=1\;}$
\texttt{EI} $\xrightarrow{\;b=1\;}$ \texttt{EP}
\end{center}
\texttt{drtran} generalizes the cast to a \emph{directed acyclic graph} of
transfers, declared in a file: each link subtracts its transfer from \emph{its}
output, so a series may be both a target and a source; the VARMA stays diagonal;
and forecasting walks the graph in topological order, propagating the system's
$\psi$ weights,
\begin{equation}
\Psi_{ij}(B) \;=\; \delta_{ij}\,\psi_i(B)
   \;+\; \sum_{k:\,\mathrm{out}(k)=i} \nu_k(B)\,\Psi_{\mathrm{in}(k),\,j}(B),
\end{equation}
so that the forecast error of a series inherits the innovations of everything
upstream of it, each filtered by the $\nu(B)$ it crosses. A \emph{cycle} is
rejected: the system would be simultaneous and cannot be triangularized by
subtraction.

Measured on a synthetic chain $X \rightarrow M \rightarrow Y$: the network
recovers both transfers and, \textbf{with the same seven free parameters}, beats
the star by $213$ points of log-likelihood ($-1038.6$ vs $-1251.8$). In the star,
the direct link $X\rightarrow Y$ that it is forced to postulate is insignificant
($t=1.45$)---because the effect of $X$ on $Y$ is \emph{indirect}, and only the
network can say so.

\subsection{Restrictions: transform the data, or declare them}

The dissertation imposes long-run neutrality ($g = \nu(1) = 1$, unit long-run
gain) by \emph{transforming the data}, rewriting the model as
\[
Y_t - \ln M1_{t-b} \;=\; \frac{\omega^{*}(B)}{\delta(B)}\,B^{b}\,
\nabla \ln M1_t \;+\; N_t,
\qquad \omega^{*}(B)=\omega(B)-\delta(B),
\]
an elegant solution of its time: the restriction is nonlinear in the parameters,
so it is taken \emph{out} of the optimizer by hand. It has the further virtue of
giving the resulting output an economic reading (real balances).

\texttt{drtran} instead declares restrictions in a \emph{slot table}: every
structural parameter is FREE, FIXED or ALIAS (shared), stated in a file using the
names the program prints. This covers the \emph{linear} restrictions---sharing
and fixing, which are what the dissertation uses for identity-linked
triads---but it does \textbf{not} cover the gain, which is nonlinear. See
\S\ref{sec:todo}.

\subsection{The start of the sample: \textsc{taste} did it better}

The transfer $\sum_k \nu_k X_{t-k}$ needs the input's past, and at $t=1$ there is
none. \textsc{taste} \emph{backforecasts} it, which is Box and Jenkins's
solution. \texttt{drtran} \emph{truncates}: it assumes a zero pre-sample input.

Measured: with $r=1$ and $\delta=0.6$ the error at $t=1$ is $63\%$ of
$\mathrm{sd}(w_X)$ and decays as $\delta^t$, contaminating the first nine or so
observations. With 215--400 data points this is irrelevant. With the 69
observations of \texttt{m6} it would be $15\%$ of the sample. \textbf{Here the
1979 solution is better than ours}, and it is on the list.

\section{An identification hazard}
\label{sec:identification}

This is the one result in the note, and it came out of an experiment that was
about something else.

Estimating Spanish CPI $\leftarrow$ WTI with a \emph{contemporaneous} transfer
($b=0$) \emph{and} freeing the innovation covariance at the same time sends the
fit to an absurd corner:

\begin{center}
\begin{tabular}{lrr}
\toprule
 & $b=0$, $\Sigma$ diagonal & $b=0$, $\sigma_{12}$ free \\
\midrule
$\log L$                  & $-718.184$ & $-718.170$ \\
$\omega_0$                & $0.0164$   & $\mathbf{0.1520}$ \\
$\mathrm{corr}(a_N,a_X)$  & $0$        & $\mathbf{-0.985}$ \\
$t$ of $\sigma_{12}$      & ---        & $\mathbf{-2424}$ \\
\bottomrule
\end{tabular}
\end{center}

The likelihood \emph{does not improve} ($\mathrm{LR}=0.029$, against
$\chi^2_1 = 3.84$), yet the estimates fly. This is not an optimizer failure. It is
a nearly flat ridge, and it has an exact explanation.

\paragraph{The prewhitened cross-correlation.} Filter both series by the input's
model: $\alpha_t = \phi_X(B) X_t = a_{X,t}$ and $\beta_t = \phi_X(B) Y_t$.

\begin{description}
\item[(A) Contemporaneous transfer] $Y_t = \omega_0 X_t + N_t$, with $N \perp X$:
\[
\beta_t = \omega_0 a_{X,t} + \phi_X(B) N_t,
\qquad
\mathrm{Cov}(\beta_t,\alpha_{t-k}) =
\begin{cases}
\omega_0\,\sigma_X^2, & k=0,\\[2pt]
0, & k>0 .
\end{cases}
\]
A clean spike at $k=0$.

\item[(B) Innovation covariance] $Y_t = N_t$, with
$\mathrm{cov}(a_{N,t}, a_{X,t}) = \sigma_{12}$:
\[
\beta_t = \phi_X(B)\,\psi_N(B)\,a_{N,t},
\qquad
\mathrm{Cov}(\beta_t,\alpha_{t-k})
 = \sigma_{12}\,\bigl[\,B^k\ \text{coefficient of}\ \phi_X(B)\psi_N(B)\,\bigr].
\]
With AR(1) noise and input this is $\sigma_{12}$ at $k=0$ and
\begin{equation}
\boxed{\;
\mathrm{Cov}(\beta_t,\alpha_{t-k}) \;=\;
\sigma_{12}\;\phi_N^{\,k-1}\,\bigl(\phi_N - \phi_X\bigr),
\qquad k \ge 1 .\;}
\label{eq:tail}
\end{equation}
The \emph{same} spike at $k=0$, plus a tail proportional to
$(\phi_N - \phi_X)$.
\end{description}

\paragraph{Consequence.} If $\phi_N = \phi_X$ the tail in \eqref{eq:tail}
vanishes and the two models are \textbf{exactly observationally equivalent}. This
is not numerical collinearity: it is non-identification. And the closer the two
autoregressive structures, the weaker the identification. The tail relative to
the spike:

\begin{center}
\begin{tabular}{lrrrrr}
\toprule
case & $\phi_N$ & $\phi_X$ & $k=1$ & $k=2$ & $k=3$ \\
\midrule
CPI $\leftarrow$ WTI (real) & $0.403$ & $0.299$ & $\mathbf{+0.104}$ & $+0.042$ & $+0.017$ \\
Synthetic                   & $0.300$ & $0.500$ & $-0.200$ & $-0.060$ & $-0.018$ \\
$\phi_N = \phi_X$           & ---     & ---     & $\mathbf{0.000}$ & $0.000$ & $0.000$ \\
\bottomrule
\end{tabular}
\end{center}

In the real case, \emph{all} of the information separating the two mechanisms is
a signal worth $10.4\%$ of the main one. With 215 observations it cannot be
resolved. Hence the ridge.

\paragraph{Three practical consequences,} all verified:
\begin{enumerate}
\item With $b \ge 1$ the problem disappears: the transfer contributes nothing at
  $k=0$, so $\sigma_{12}$ is cleanly identified. On synthetic data with $b=2$ and
  a true $\sigma_{12}=0$: $t=1.33$, $p=0.18$, correctly insignificant, with the
  $\omega$'s untouched.
\item \texttt{drtran} now \emph{warns} when a link with $b=0$ has its covariance
  freed.
\item It explains the school's practice. \texttt{m6-1} carries off-diagonal
  covariances and \emph{no} contemporaneous structure whatever; and the
  dissertation~\cite[\S2.4]{munoz2001} says that the specification of a bivariate
  relation \emph{``may begin with the modification of the matrix $\Sigma$''}. One
  uses \emph{one} of the two devices, never both. We rediscovered this by walking
  into it.
\end{enumerate}

\section{An application: the pass-through of crude oil to Spanish CPI is not rational}

Three fits, same output and same input, the first two with \emph{exactly} the same
number of free parameters:

\begin{center}
\begin{tabular}{lrrl}
\toprule
model & free & $\log L$ & adequacy \\
\midrule
$\omega_0/(1-\delta_1 B)$ \quad (rational)      & 17 & $-721.72$ & \textbf{INADEQUATE} ($p=0.033$) \\
$\omega_0 + \omega_1 B$ \quad (two $\omega$'s)  & 17 & $\mathbf{-718.18}$ & adequate ($p=0.196$) \\
$(\omega_0+\omega_1 B)/(1-\delta_1 B)$ \quad (nests the above) & 18 & $-718.18$ & adequate \\
\bottomrule
\end{tabular}
\end{center}

The rational form \emph{loses with the same number of parameters}. The third fit,
which nests the second, settles it: $\delta_1 = -0.010$ with $t=-0.055$;
$\mathrm{LR}=0.003$ against $\chi^2_1=3.84$. There is no denominator.

The reason is visible in the impulse weights:
\begin{center}
\begin{tabular}{lrrrrr}
\toprule
 & $\nu_0$ & $\nu_1$ & $\nu_2$ & $\nu_3$ & $\nu_4$\\
\midrule
rational        & $0.01680$ & $0.00769$ & $0.00352$ & $0.00161$ & $0.00074$\\
two $\omega$'s  & $0.01640$ & $0.01079$ & $0$ & $0$ & $0$\\
\bottomrule
\end{tabular}
\end{center}
The data want $\nu_1/\nu_0 \approx 0.66$ \emph{and} $\nu_2 = 0$. A geometric decay
cannot do both. The adequacy test catches it exactly there: the noise still
carries a trace of the input at $k=2$, with $r=-0.14$---negative, because the
rational form puts \emph{too much} weight at that lag.

\paragraph{A methodological warning.} In the \emph{forced} rational fit,
$\delta_1 = 0.458$ with $t=6.40$: apparently overwhelming. It is a mirage. With
the denominator imposed, $\delta$ is the \emph{only} route to any weight at lag
one; its significance says ``there is a response at lag one'', not ``the response
is geometric''. The honest test is the one in the nesting model, and there it
dies.

Economically: the pass-through of crude to the Spanish CPI is immediate and
exhausted within two months---not a distributed lag that propagates. The estimate
$\omega = (0.0164,\,0.0108)$ reproduces the previously documented values
$(0.0154,\,0.0104)$.

\section{What the dissertation teaches and the program still does not do}
\label{sec:todo}

\begin{enumerate}
\item \textbf{The order of reformulation.} The dissertation~\cite[\S2.6]{munoz2001}
  states an asymmetry that is not obvious: a misspecified \emph{relation} can make
  the residual ACF/PACF look like a misspecified \emph{noise}, but \emph{``a
  misspecified noise cannot give, in the CCF, the impression of a misspecified
  relation. For these reasons, $\nu(B)$ is reformulated until it appears adequate
  \emph{before} reformulating $\theta(B)$.''} The program prints both diagnostics
  and says nothing about which to act on first. It should: \emph{the relation
  first, the noise afterwards}.

\item \textbf{The long-run gain $g=\nu(1)$} is not reported. It is the quantity
  with economic meaning---the total permanent effect of the input on the
  output---and the one the dissertation wants to test. It is
  $\sum\omega / (1-\sum\delta)$, and its standard error follows from the delta
  method with the covariance matrix already computed. With it, a long-run
  neutrality hypothesis can be tested \emph{directly}, without transforming the
  data and without estimating two models.

\item \textbf{Identity-linked triads.} The dissertation handles them
  (\S2.4) by restricting the deterministic parameters of the three series to
  satisfy the identity, \emph{``economizing the number of parameters to be
  estimated''}. The slot table can already express this; it has not been
  exercised, and the linear combinations of forecasts with their $c'Vc$ variance
  (the aggregates) are not implemented.
\end{enumerate}

\section*{Reproducibility}

Every figure in this note is produced by the program's test battery
(\texttt{test\_battery.sh}, 134 checks): homologation against the univariate
program (\S1--1d), synthetic ground truth (\S2--2c), multiple inputs (\S2d),
shared and fixed parameters (\S2e), the network (\S2f), the covariance (\S2g),
the real case and its rational form (\S3b), separability (\S3c), and the
identification hazard (\S3d). The source is at
\url{https://github.com/davidesg/drtran}, under the GNU GPL.

\begin{thebibliography}{9}

\bibitem{bjr1994}
G.~E.~P.~Box, G.~M.~Jenkins and G.~C.~Reinsel (1994).
\emph{Time Series Analysis: Forecasting and Control}, 3rd ed. Prentice Hall.
Chapters 10--11.

\bibitem{dennis1983}
J.~E.~Dennis and R.~B.~Schnabel (1983).
\emph{Numerical Methods for Unconstrained Optimization and Nonlinear Equations}.
Prentice Hall.

\bibitem{mauricio1995}
J.~A.~Mauricio (1995).
Exact maximum likelihood estimation of stationary vector ARMA models.
\emph{Journal of the American Statistical Association} \textbf{90}, 282--291.
\emph{(The algorithm implemented in the likelihood engine used here.)}

\bibitem{mauricio2002}
J.~A.~Mauricio (2002).
An algorithm for the exact likelihood of a stationary vector
autoregressive-moving average model.
\emph{Journal of Time Series Analysis} \textbf{23}(4), 473--486.
\emph{(An innovations-form algorithm; a different one, not the one wired in.)}

\bibitem{melard1984}
G.~M\'elard (1984).
Algorithm AS 197: A fast algorithm for the exact likelihood of
autoregressive-moving average models.
\emph{Applied Statistics} \textbf{33}, 104--114.
\emph{(The univariate engine behind \texttt{fue}.)}

\bibitem{munoz2001}
M.~S.~Mu\~noz Polo (2001).
\emph{Estudios econom\'etricos de las series temporales de la industria
espa\~nola}. Doctoral dissertation, Universidad Complutense de Madrid.
Supervisor: A.~B.~Treadway.

\bibitem{tiaobox1981}
G.~C.~Tiao and G.~E.~P.~Box (1981).
Modeling multiple time series with applications.
\emph{Journal of the American Statistical Association} \textbf{76}, 802--816.

\end{thebibliography}

\end{document}
